\section{Coverage Definitions and Native Source Audit}
\label{app:systems}

% Role: provenance. Message: the comparison is fixed to inspected source revisions.
Table~\ref{tab:systems} uses a source audit dated November 2, 2026. Penguin Harness is pinned to \path{7406d13e2f4b3ba0fe495feed32c172f3484be3b}, whose core package declares version 0.2.18. Proteus uses the report snapshot in Appendix~\ref{app:evidence}. The source bundle includes file hashes, cell definitions, and links. No competing framework was executed and no task scores were compared in this audit.

% Role: definition. Message: the symbols express narrow propositions rather than overall rank.
An infrastructure check means that the inspected source or documentation supplies the named mechanism; it does not establish equal enforcement, reliability, or coverage. A private-evaluation check denotes an information protocol, not a security guarantee. A method cross means absence from the inspected built-in plugin library. A question mark means that the precise common service was not established, not absence of all related workflows, logging, budget limits, or session recovery. Both systems are extensible; the table's five-method subset is not an overall breadth ranking.

\paragraph{Table sources and scope.}
The comparison uses the pinned architecture and self-improvement guides, built-in toolkit library, and trace-resume implementation~\citep{penguin_software}. The source manifest preserves the inspected files and each cell's definition. Proteus's corresponding stage, workflow, evidence, and operation contracts are described in Section~\ref{sec:framework}. A table check establishes a declared mechanism, not a stronger isolation guarantee or a controlled quality result.

\paragraph{Native algorithm sources.}
The reference protocol is defined in \repo{proteus/core/episode_protocol.py}{\code{episode\_protocol.py}}; goal/evidence injection is implemented in \repo{proteus/core/episode.py}{\code{episode.py}}. The common transition and its constraints are in \repo{proteus/methods/native_evolution.py}{\code{native\_evolution.py}}, with active/pending workspace execution in \repo{proteus/execution/harness_evolution.py}{\code{harness\_evolution.py}}. The \repo{recipes/native/shared-evaluation/README.md}{Native shared-evaluation fixture} provides a model-free invalid-edit, repair, and private-grading example. It is mechanism evidence, not learned-quality evidence. The common profile's observed-evaluation requirement, single episode stage, and last-valid output must remain distinct from legacy Native's broader configuration options.

\paragraph{Algorithm versus system contribution.}
Native's research proposition concerns goal interpretation, evidence use, self-authored tests, and viability-gated continuation by the evolving harness. The framework makes these choices inspectable and supplies other methods for comparison. A reference method can in principle be integrated into another extensible toolkit; its scientific value therefore rests on a precise protocol and controlled evidence, not exclusivity to one software package. The four phase names and source activation plumbing alone do not establish algorithmic novelty.

\paragraph{Comparison still needed.}
Use identical starting artifacts, tasks, models, and budgets to vary broad-goal versus score-target wording, last-valid versus score-based selection, and self-evidence availability. Preserve the independent hidden evaluator and report regressions as well as improvements. Integration-effort and post-dispatch interruption studies would separately test the framework's composition and recovery claims. None of these proposed comparisons is a result of the present report.
